\chapter{Đối chiếu đặc trưng và nguồn dữ liệu}
\label{app:features}

\section{Danh sách 23 đặc trưng đầu vào ban đầu}
Bảng~\ref{tab:features} dùng cùng danh sách \code{FEATURE_COLUMNS} trong notebook. Nguồn được ghi theo bảng gốc để không trộn tên bảng trung gian với tên dữ liệu đầu vào. Các thuộc tính đã chọn chưa được mã hóa hoặc lựa chọn tối ưu bằng thực nghiệm.

\begin{longtable}{@{}p{.25\linewidth}p{.20\linewidth}p{.49\linewidth}@{}}
\caption{Nguồn và định nghĩa bộ đặc trưng đầu vào ứng viên.}\label{tab:features}\\
\toprule
\textbf{Đặc trưng} & \textbf{Nguồn gốc} & \textbf{Loại hình thành và định nghĩa} \\
\midrule
\endfirsthead
\multicolumn{3}{l}{\small Bảng~\ref{tab:features} (tiếp theo)}\\
\toprule
\textbf{Đặc trưng} & \textbf{Nguồn gốc} & \textbf{Loại hình thành và định nghĩa} \\
\midrule
\endhead
\midrule
\multicolumn{3}{r}{\small Còn tiếp ở trang sau}\\
\endfoot
\bottomrule
\endlastfoot
\code{gender} & \code{studentInfo} & Thuộc tính gốc. Giới tính \\
\code{region} & \code{studentInfo} & Thuộc tính gốc. Vùng cư trú \\
\code{highest_education} & \code{studentInfo} & Thuộc tính gốc. Trình độ học vấn cao nhất \\
\code{imd_band} & \code{studentInfo} & Chuẩn hóa nhãn. Unknown cho thiếu; nhóm IMD thống nhất hậu tố \% \\
\code{age_band} & \code{studentInfo} & Thuộc tính gốc. Nhóm tuổi \\
\code{num_of_prev_attempts} & \code{studentInfo} & Thuộc tính gốc. Số lần học trước \\
\code{studied_credits} & \code{studentInfo} & Thuộc tính gốc. Số tín chỉ đang học \\
\code{disability} & \code{studentInfo} & Thuộc tính gốc. Khai báo khuyết tật \\
\code{imd_band_missing} & \code{studentInfo} & Cờ mới. Cờ thiếu IMD ban đầu; gán cố định, không học từ y \\
\code{date_registration} & \code{studentRegistration} & Thuộc tính gốc. Ngày đăng ký đã biết và không sau ngày chốt \\
\code{module_presentation_length} & \code{courses} & Thuộc tính gốc. Thời lượng dự kiến của đợt học \\
\code{n_submissions} & \code{studentAssessment} + \code{assessments} & Tổng hợp mới. Số bài TMA/CMA không banked nộp trong ngày 0 đến ngày chốt \\
\code{n_submissions_with_due_date} & \code{studentAssessment} + \code{assessments} & Tổng hợp mới. Số bài nộp trên có hạn xác định và hạn không sau ngày chốt \\
\code{n_late_submissions} & \code{studentAssessment} + \code{assessments} & Tổng hợp mới. Số bài có ngày nộp > hạn, trong nhóm có hạn đến ngày chốt \\
\code{late_submission_rate} & \code{studentAssessment} + \code{assessments} & Tổng hợp mới. Số bài muộn / số bài có hạn đến ngày chốt; để thiếu nếu mẫu số bằng 0 \\
\code{total_clicks} & \code{studentVle} + \code{vle} & Tổng hợp mới. Tổng click ghi nhận trong ngày 0 đến ngày chốt, giữ log trùng chưa xác minh \\
\code{active_days} & \code{studentVle} + \code{vle} & Tổng hợp mới. Số ngày khác nhau có bản ghi VLE trong ngày 0 đến ngày chốt \\
\code{resource_types_used} & \code{studentVle} + \code{vle} & Tổng hợp mới. Số activity\_type khác nhau đã có tương tác trong cửa sổ \\
\code{avg_clicks_per_day} & \code{studentVle} + \code{vle} & Tổng hợp mới. total\_clicks / (ngày chốt+1), gồm ngày không có log \\
\code{clicks_oucontent} & \code{studentVle} + \code{vle} & Tổng hợp mới. Tổng click của loại oucontent trong ngày 0 đến ngày chốt \\
\code{clicks_forumng} & \code{studentVle} + \code{vle} & Tổng hợp mới. Tổng click của loại forumng trong ngày 0 đến ngày chốt \\
\code{clicks_quiz} & \code{studentVle} + \code{vle} & Tổng hợp mới. Tổng click của loại quiz trong ngày 0 đến ngày chốt \\
\code{clicks_resource} & \code{studentVle} + \code{vle} & Tổng hợp mới. Tổng click của loại resource trong ngày 0 đến ngày chốt \\

\end{longtable}

\section{Đường dẫn và khả năng đối chiếu}
Nguồn số liệu chính là \code{preprocess.ipynb}, các CSV trong \code{report/tables/}, hình trong \code{report/figures/} và năm tệp theo tuần trong \code{data/preprocessed/}. Tên trường kỹ thuật trong dữ liệu mô hình được giữ nguyên; tên cột báo cáo được trình bày bằng tiếng Việt.

Các bảng chọn lọc và giá trị định lượng trong báo cáo được tạo lại bằng \code{scripts/prepare_tables.py}. Tệp \code{source_manifest.json} lưu danh sách nguồn và SHA-256 để đối chiếu phiên bản. Hướng dẫn biên dịch, danh sách hình/bảng sử dụng và các giới hạn xác minh được ghi trong \code{README.md}, \code{ASSETS.md} và \code{output/validation.json} của bộ mã nguồn báo cáo.

Mã sinh viên trong dữ liệu chỉ là mã định danh đã ẩn danh. Thông tin tác giả và giảng viên trên bìa cần được người thực hiện điền trong \code{metadata.tex}; không lấy từ báo cáo tham khảo.
